% GENERATED FILE. Edit canonical lessons, then run npm run generate:books.
% canonical-source-hash: fnv1a64:af6ee37dbb1d7fab
% canonical-lessons: HI-C448-conc, HI-C448-pankh, HI-C448-sikar, HI-C448-cugna, HI-C448-cubhna

\chapter{Beaks, Wings and Thorns}
\label{ch:hi-beaks-wings-and-thorns}
\hlchaptermodality{448}

\begin{chapteropening}
\textbf{By the end of this chapter:} \emph{I can talk about beaks and wings, a hunt, birds pecking grain and a thorn that pricks.}

Complete the last lesson of chapter 448: I can talk about beaks and wings, a hunt, birds pecking grain and a thorn that pricks.
\end{chapteropening}

\section[\texorpdfstring{\hi{चोंच} (coñc)}{चोंच (coñc)}]{\hi{चोंच} (coñc) — a beak}

\label{lesson:HI-C448-conc}



\begin{quote}
Before the new one: say the Hindi for an insect, then the Hindi for to wind round.
\end{quote}

\subsection*{You'll want to know: \hi{चोंच}}
\textbf{\hi{चोंच}} — \emph{coñc} — \textquotedblleft{}a beak\textquotedblright{}.

\begin{quote}
\hi{तोते} \hi{की} \hi{चोंच} \hi{लाल} \hi{होती} \hi{है।}
\end{quote}

\begin{quote}
\emph{tote kī coñc lāl hotī hai.}
\end{quote}

\begin{quote}
A parrot's beak is red.
\end{quote}

Feminine.

\begin{cousinweb}
\textbf{\hi{चोंच}} (\emph{coñc}) is a tadbhava of Sanskrit \emph{cañcu}, \textquotedblleft{}a beak\textquotedblright{} (CDIAL 4569), which lists Hindi \emph{cõc}; CDIAL thinks the Sanskrit word itself came from Dravidian.
\end{cousinweb}

\begin{grammarlens}[title={What you've built}]
One more word for getting around.
\end{grammarlens}

\subsection*{Your turn}
\begin{itemize}
\raggedright
  \item \emph{Say it:} \emph{coñc}
  \item \emph{Say it:} \emph{coñc}, once more
  \item \emph{Say it:} \emph{tote kī coñc lāl hotī hai.} — A parrot's beak is red.
  \item \emph{Recall:} say the Hindi for an insect, then the Hindi for to wind round, then say \emph{coñc} again
\end{itemize}

\begin{tcolorbox}[breakable,colback=teal!4,colframe=teal!35!black,title={Before you move on}]
What does \hi{चोंच} mean? (\textquotedblleft{}A beak\textquotedblright{}.) Say it once more.
\end{tcolorbox}
\section[\texorpdfstring{\hi{पंख} (paṅkh)}{पंख (paṅkh)}]{\hi{पंख} (paṅkh) — a wing; a feather}

\label{lesson:HI-C448-pankh}



\begin{quote}
Before the new one: say the Hindi for to wind round, then the Hindi for a beak.
\end{quote}

\subsection*{You'll want to know: \hi{पंख}}
\textbf{\hi{पंख}} — \emph{paṅkh} — \textquotedblleft{}a wing; a feather\textquotedblright{}.

\begin{quote}
\hi{चिड़िया} \hi{ने} \hi{अपने} \hi{पंख} \hi{फैलाए।}
\end{quote}

\begin{quote}
\emph{ciṛiyā ne apne paṅkh phailāe.}
\end{quote}

\begin{quote}
The bird spread its wings.
\end{quote}

Masculine.

\begin{cousinweb}
\textbf{\hi{पंख}} (\emph{paṅkh}) comes from Sanskrit \emph{pakṣa}, \textquotedblleft{}a wing, a feather, a side\textquotedblright{} (CDIAL 7627), which lists Hindi \emph{pā̃kh}, \textquotedblleft{}wing\textquotedblright{}, and \hi{पंखा}, a fan: the fan is a wing that moves the air.
\end{cousinweb}

\begin{grammarlens}[title={What you've built}]
One more word for getting around.
\end{grammarlens}

\subsection*{Your turn}
\begin{itemize}
\raggedright
  \item \emph{Say it:} \emph{paṅkh}
  \item \emph{Say it:} \emph{paṅkh}, once more
  \item \emph{Say it:} \emph{ciṛiyā ne apne paṅkh phailāe.} — The bird spread its wings.
  \item \emph{Recall:} say the Hindi for to wind round, then the Hindi for a beak, then say \emph{paṅkh} again
\end{itemize}

\begin{tcolorbox}[breakable,colback=teal!4,colframe=teal!35!black,title={Before you move on}]
What does \hi{पंख} mean? (\textquotedblleft{}A wing; a feather\textquotedblright{}.) Say it once more.
\end{tcolorbox}
\section[\texorpdfstring{\hi{शिकार} (śikār)}{शिकार (śikār)}]{\hi{शिकार} (śikār) — hunting; prey}

\label{lesson:HI-C448-sikar}



\begin{quote}
Before the new one: say the Hindi for a beak, then the Hindi for a wing.
\end{quote}

\subsection*{You'll want to know: \hi{शिकार}}
\textbf{\hi{शिकार}} — \emph{śikār} — \textquotedblleft{}hunting; prey\textquotedblright{}.

\begin{quote}
\hi{शेर} \hi{ने} \hi{एक} \hi{हिरन} \hi{का} \hi{शिकार} \hi{किया।}
\end{quote}

\begin{quote}
\emph{śer ne ek hiran kā śikār kiyā.}
\end{quote}

\begin{quote}
The lion hunted down a deer.
\end{quote}

Masculine.

\begin{cousinweb}
\textbf{\hi{शिकार}} (\emph{śikār}) is usually traced to Persian \emph{šikār}, \textquotedblleft{}hunting, the hunt, prey\textquotedblright{}.
\end{cousinweb}

\begin{grammarlens}[title={What you've built}]
One more word for getting around.
\end{grammarlens}

\subsection*{Your turn}
\begin{itemize}
\raggedright
  \item \emph{Say it:} \emph{śikār}
  \item \emph{Say it:} \emph{śikār}, once more
  \item \emph{Say it:} \emph{śer ne ek hiran kā śikār kiyā.} — The lion hunted down a deer.
  \item \emph{Recall:} say the Hindi for a beak, then the Hindi for a wing, then say \emph{śikār} again
\end{itemize}

\begin{tcolorbox}[breakable,colback=teal!4,colframe=teal!35!black,title={Before you move on}]
What does \hi{शिकार} mean? (\textquotedblleft{}Hunting; prey\textquotedblright{}.) Say it once more.
\end{tcolorbox}
\section[\texorpdfstring{\hi{चुगना} (cugnā)}{चुगना (cugnā)}]{\hi{चुगना} (cugnā) — to peck at; to pick up (grain)}

\label{lesson:HI-C448-cugna}



\begin{quote}
Before the new one: say the Hindi for a wing, then the Hindi for hunting.
\end{quote}

\subsection*{You'll want to know: \hi{चुगना}}
\textbf{\hi{चुगना}} — \emph{cugnā} — \textquotedblleft{}to peck at; to pick up (grain)\textquotedblright{}.

\begin{quote}
\hi{आँगन} \hi{में} \hi{कबूतर} \hi{चावल} \hi{चुग} \hi{रहे} \hi{थे।}
\end{quote}

\begin{quote}
\emph{āṁgan meṁ kabūtar cāval cug rahe the.}
\end{quote}

\begin{quote}
Pigeons were pecking at rice in the courtyard.
\end{quote}

Past (it …): \hi{वह} \hi{चिड़िया} \hi{पूरा} \hi{अनाज} \hi{चुग} \hi{गई}, \emph{vah ciṛiyā pūrā anāj cug gaī}.

\begin{cousinweb}
\textbf{\hi{चुगना}} (\emph{cugnā}) goes back to \emph{*cugyati}, \textquotedblleft{}pecks\textquotedblright{}, a form CDIAL reconstructs (CDIAL 4852), which lists Hindi \emph{cugnā}, \textquotedblleft{}to peck, pick up, feed\textquotedblright{}.
\end{cousinweb}

\begin{grammarlens}[title={What you've built}]
One more word for getting around.
\end{grammarlens}

\subsection*{Your turn}
\begin{itemize}
\raggedright
  \item \emph{Say it:} \emph{cugnā}
  \item \emph{Say it:} \emph{cugnā}, once more
  \item \emph{Say it:} \emph{āṁgan meṁ kabūtar cāval cug rahe the.} — Pigeons were pecking at rice in the courtyard.
  \item \emph{Recall:} say the Hindi for a wing, then the Hindi for hunting, then say \emph{cugnā} again
\end{itemize}

\begin{tcolorbox}[breakable,colback=teal!4,colframe=teal!35!black,title={Before you move on}]
What does \hi{चुगना} mean? (\textquotedblleft{}To peck at; to pick up (grain)\textquotedblright{}.) Say it once more.
\end{tcolorbox}
\section[\texorpdfstring{\hi{चुभना} (cubhnā)}{चुभना (cubhnā)}]{\hi{चुभना} (cubhnā) — to prick; to stick into}

\label{lesson:HI-C448-cubhna}



\begin{quote}
Before the new one: say the Hindi for hunting, then the Hindi for to peck.
\end{quote}

\subsection*{You'll want to know: \hi{चुभना}}
\textbf{\hi{चुभना}} — \emph{cubhnā} — \textquotedblleft{}to prick; to stick into\textquotedblright{}.

\begin{quote}
\hi{एक} \hi{काँटा} \hi{मेरे} \hi{पैर} \hi{में} \hi{चुभ} \hi{गया।}
\end{quote}

\begin{quote}
\emph{ek kā̃ṭā mere pair meṁ cubh gayā.}
\end{quote}

\begin{quote}
A thorn pricked my foot.
\end{quote}

Past (it …): \hi{वह} \hi{काँटा} \hi{उँगली} \hi{में} \hi{चुभा}, \emph{vah kā̃ṭā uṅglī meṁ cubhā}.

\begin{cousinweb}
\textbf{\hi{चुभना}} (\emph{cubhnā}) goes back to \emph{*cubhyate}, \textquotedblleft{}is pierced\textquotedblright{}, a form CDIAL reconstructs (CDIAL 4867), which lists Hindi \emph{cubhnā}. A harsh word \hi{चुभता} too: it stings.
\end{cousinweb}

\begin{grammarlens}[title={What you've built}]
One more word for getting around.
\end{grammarlens}

\subsection*{Your turn}
\begin{itemize}
\raggedright
  \item \emph{Say it:} \emph{cubhnā}
  \item \emph{Say it:} \emph{cubhnā}, once more
  \item \emph{Say it:} \emph{ek kā̃ṭā mere pair meṁ cubh gayā.} — A thorn pricked my foot.
  \item \emph{Recall:} say the Hindi for hunting, then the Hindi for to peck, then say \emph{cubhnā} again
\end{itemize}

\begin{tcolorbox}[breakable,colback=teal!4,colframe=teal!35!black,title={Before you move on}]
What does \hi{चुभना} mean? (\textquotedblleft{}To prick; to stick into\textquotedblright{}.) Say it once more.
\end{tcolorbox}
